% 单位容量撤回例：先 s->a->b->t，再沿残量路径 s->b->a->t（经反向残量 b->a），流值 2（Python 验算）。
\begin{tikzpicture}[>=Stealth,
  v/.style={circle,draw=FigureStroke,fill=FigureBlueFill,line width=.8pt,minimum size=6mm,inner sep=0pt},
  arc/.style={->,draw=FigureGrid,line width=.8pt},
  flow/.style={->,draw=FigureBlue,line width=1.6pt},
  res/.style={->,draw=FigureMuted,dashed,line width=.8pt},
  path/.style={->,draw=FigureGreen,line width=1.6pt},
  w/.style={fill=white,inner sep=.6pt,font=\figurefont\fontsize{7}{8}\selectfont},
  lab/.style={font=\figurefont\fontsize{7.5}{9}\selectfont}]
\newcommand{\verts}{%
  \node[v] (s) at (0,0) {$s$}; \node[v] (a) at (1.4,.95) {$a$};
  \node[v] (b) at (1.4,-.95) {$b$}; \node[v] (t) at (2.8,0) {$t$};}
\begin{scope}
  \verts
  \draw[flow] (s)--node[w]{$1/1$}(a); \draw[flow] (a)--node[w]{$1/1$}(b); \draw[flow] (b)--node[w]{$1/1$}(t);
  \draw[arc] (s)--node[w]{$0/1$}(b); \draw[arc] (a)--node[w]{$0/1$}(t);
  \node[lab,align=center] at (1.4,-1.85) {(a) 第一次：$s\to a\to b\to t$\\流值$1$};
\end{scope}
\begin{scope}[xshift=3.9cm]
  \verts
  \draw[res] (a)--node[w]{$1$}(s); \draw[res] (t)--node[w]{$1$}(b);
  \draw[path] (s)--node[w]{$1$}(b); \draw[path] (b)--node[w]{$1$}(a); \draw[path] (a)--node[w]{$1$}(t);
  \node[lab,align=center] at (1.4,-1.85) {(b) 残量网络中的\\路径$s\to b\to a\to t$};
\end{scope}
\begin{scope}[xshift=7.8cm]
  \verts
  \draw[flow] (s)--node[w]{$1/1$}(a); \draw[flow] (a)--node[w]{$1/1$}(t);
  \draw[flow] (s)--node[w]{$1/1$}(b); \draw[flow] (b)--node[w]{$1/1$}(t);
  \draw[arc] (a)--node[w]{$0/1$}(b);
  \node[lab,align=center] at (1.4,-1.85) {(c) 撤回$a\to b$上的流量\\流值$2$};
\end{scope}
\end{tikzpicture}
